\section{The finite induction from five to fourteen}
\label{ext:finite-five-fourteen}
The fifth index starts a different part of the proof. The polynomial
majorants for two and three summands and the physical comparison at four
are followed by the conditional count comparison of
Section~\ref{ext:finite-counts}.
The comparison itself is the same at every index; what changes is the
numerical region in which the law has first been placed. This section
describes how those regions are acquired and how they cover all third
moments.

\subsection{Closing the fifth index}
\label{ext:five-assembly}
The fifth-index calculation has three complementary parts: a count
comparison at small third moments, an elementary corner estimate where
the signed moment is large and positive, and a Fourier estimate for the
remaining third moments. The corner is a genuine part of the cover;
it cannot be discarded as a small exceptional region.

\begin{proposition}[Fifth-index assembly criterion]
\label{ext:five-assembly-criterion}
Assume $C_4\le\CE$, the continuous inequalities proved in this paper,
and complete verification of the finite inputs and their implications
described below. Then $C_5\le\CE$.
\end{proposition}
\begin{proof}
Suppose $C_5>\CE$ and choose a positive closed-event ratio maximizer.
Its same-law predecessor score is at most $C_4$, so its contact
allowance is zero. The support calculations in
Appendix~\ref{ext:n5-support} therefore apply to this very marginal.
They precede the supported Fourier estimates. The unrestricted
auxiliary Fourier estimate first supplies $C_5<429/1000$; this is
the input to the two lowest-moment initializers, not their conclusion.
The higher-moment initializers use separately acquired supported
own-score bounds. Their scalar input ceilings are
\[
\begin{array}{c|cc}
\text{region}&\text{own score}&\text{HIGH}\\\hline
1\le\rho\le11/10&429/1000&337/1000\\
11/10\le\rho\le23/20&429/1000&369/1000\\
23/20\le\rho\le237/200&52/125&393/1000\\
237/200\le\rho\le6/5,\ \beta\le1/5&207/500&193/500\\
237/200\le\rho\le6/5,\ \beta\ge1/5&103/250&193/500
\end{array}
\]
Here HIGH is the compensated high-frequency contribution, as in
Lemma~\ref{ext:compensated-high}. Each such input is used only on a
support and moment domain contained in the domain of its Fourier
proof. In particular a better score bound requiring a smaller support
radius is not transferred to a larger-radius profile.
The last two rows are restricted to the count rectangles specified below;
the separate corner does not require these inputs.

Here is the exact signed-moment partition behind the small-moment
count calculation. Set
\[
 b_2=\frac{84729828287999}{250000000000000},\qquad
 b_3=\frac{35139006743001}{10^{14}}.
\]
The sharp signed-moment inequality bounds $|\beta|$ by $b_2$ on
$\rho\le237/200$ and by $b_3$ on $\rho\le6/5$.
The rational endpoints are checked by
$(1+b_i^2)^2\ge\rho_i^2(1-b_i^2)$ at the respective upper
third moments, using the monotonicity of this expression in $b_i\ge0$.

The two lowest moment intervals each form one full signed region.
On $[23/20,237/200]$, split $\rho$ at $467/400$ and $\beta$ at
$1/5$, giving four closed rectangles. On $[237/200,6/5]$, split
$\rho$ at $477/400$. In the lower half use the three signed intervals
$[-b_3,1/5]$, $[1/5,3/10]$, and $[3/10,b_3]$. In the upper half
use $[-b_3,1/5]$, $[1/5,3/10]$, and $[3/10,17/50]$.
These are the other six count regions. The remaining rectangle is
\begin{equation}\label{ext:five-corner-domain}
 [477/400,6/5]\times[17/50,b_3]
 \quad\text{in }(\rho,\beta).
\end{equation}
It is covered by the whole-line elementary estimate in
Appendix~\ref{ext:gapfree5}, including its two
unbounded threshold rays. This estimate has no support or predecessor
premise. Its upper bound is below $4097/10000<\CE$.

In each of the twelve count regions, an unseeded finite initializer
first acquires the raw remainder and then the physical defect, gap
and threshold constraints. The whole final moment cover is exported.
For each probability panel $[j/100,(j+1)/100]$, the subsequent resource
calculation concerns that same cover and that same law. Its signed
cubic, conditional variance and inward-displacement inequalities give
a joint count domain as in Section~\ref{ext:physical-count-domain}.
Every inequality in the numerical count contract must follow in its
stated direction from this acquired domain; matching a filename or
an isolated scalar maximum is insufficient.

The twelve regions have $1,200$ probability-panel instances. Of
these, $905$ are excluded by their resource inequalities. The other
$295$ are covered by six-root count forests with $1,770$ roots and
$75,178$ terminal leaves. Every accepted leaf has score strictly below
the fixed count target, and that target is strictly below $\CE$.
The interval boundaries in $\rho$, $\beta$ and $p$ are included.
Consequently these twelve regions and \eqref{ext:five-corner-domain}
cover all $1\le\rho\le6/5$, without assuming that the sign of
$\beta$ can be changed while keeping the event fixed.

For completeness, the supported Fourier complement has six closed
bands, each supplied with a previously acquired support radius:
\[
\begin{array}{c|c}
\rho& R\\\hline
{[6/5,5/4]}&13/5\\
{[5/4,7/5]}&117/40\\
{[7/5,3/2]}&25/8\\
{[3/2,8/5]}&133/40\\
{[8/5,17/10]}&7/2\\
{[17/10,7361/4000]}&15/4
\end{array}
\]
The first band uses the additional radius-$13/5$ support calculation.
The second restricts the wider support interval $[6/5,7/5]$.
All six implications retain both support signs and every real
threshold. The Fourier proof subdivides these bands into $41$ closed
strips and bounds the score below $4097/10000$ on each. Its recorded
$42$ mesh attempts include the unsuccessful coarser attempt; only
the successful complete enclosures are used as inequalities.
For $\rho\ge7361/4000$, the decreasing published affine bound is
at most $819/2000<4097/10000$.

Thus the count regions, the elementary corner, the six Fourier
bands and the infinite tail exclude the chosen maximizing witness.
Closed-event attainment and reflection account for either sign of
the original Kolmogorov discrepancy. This proves the implication.
\end{proof}

\subsection{Induction from six through fourteen}


\begin{proposition}[Finite induction, with explicit numerical premises]
\label{ext:finite-five-fourteen-assembly}
Assume $C_5\le\CE$. For each integer $n=6,\ldots,14$, assume that
the support, Fourier, initializer, resource, and count checks described
below pass on their complete domains, and that every output is joined
to the input of the next check for the same index and marginal law.
Then $C_n\le\CE$ for $6\le n\le14$.
\end{proposition}

The complete checks and their input joins are recorded in
Appendix~\ref{revision:reproducibility}. We state the assembly implication
with explicit numerical premises so that each computation can be matched
to the precise hypothesis it supplies.

\begin{proof}
Suppose there is a first failure in this range and let $n$ be its
index. Corollary~\ref{ext:ratio-attainment} supplies an attained
positive closed-event ratio maximizer $(P,H)$ with
$C=C_n>\CE$. Write $\rho=\E|X|^3$ and
$A=\sqrt{n-1}\Delta_{n-1}(P)/\rho$. The induction hypothesis gives
\[
 A\le C_{n-1}\le\CE<C,\qquad d_P=n(A-C)<0.
\]
This is the only predecessor estimate needed here. In particular it
implies the nonnegative allowances retained in the original numerical
formulas: $579/5000$ at $n=6$, and $99/1000$ or $63/1000$ in the two
low-moment bands at $n=10$. At other indices the allowance is zero.
Keeping a larger allowance weakens an upper bound; it does not introduce
an assumption about a different marginal or a later sample size.

\paragraph{The low-moment bands.}
The required data are as follows. An own-score ceiling means an upper
bound on $\sqrt n\Delta_n(P)/\rho$. A HIGH ceiling bounds the
compensated Fourier integral defined in
Lemma~\ref{ext:compensated-high}; it is not by itself a score bound.
\begin{center}\small
\begin{tabular}{@{}ccccc@{}}
\toprule
indices & $\rho$ interval & support radius & own-score ceiling & HIGH ceiling\\
\midrule
$6,\ldots,9$ & $[1,11/10]$ & $23/10$ & $423/1000$ & $69/200$\\
$6,\ldots,9$ & $[11/10,23/20]$ & $12/5$ & $209/500$ & $3/8$\\
$6,\ldots,9$ & $[23/20,237/200]$ & $5/2$ & $83/200$ & $83/200$\\
$10,\ldots,14$ & $[1,11/10]$ & $23/10$ & $1047/2500$ & $69/200$\\
$10,\ldots,14$ & $[11/10,23/20]$ & $12/5$ & $519/1250$ & $93/250$\\
\bottomrule
\end{tabular}
\end{center}
The support column is acquired first, by
Proposition~\ref{ext:support6-14-table}. The two Fourier columns
are then evaluated with the actual integer $n$ and that supported law,
using Appendix~\ref{ext:fourier-moments}. A Fourier estimate may
cover a larger moment interval or support radius than the row requires;
the join checks inclusion of the whole smaller region. In particular,
the third row may use the wider interval $[23/20,6/5]$ with radius
$13/5$ for its Fourier inputs. No estimate at index nine is transferred
numerically to index ten.

The order after these three inputs is important. First an unseeded
initializer proves a raw remainder bound, and only then feeds the
resulting defect and threshold bounds back into subsequent passes.
Appendix~\ref{ext:finite-acquisition} proves this finite enclosure
induction. The resource calculation consumes the complete initializer
and acquires a central gap, signed moments, and the common conditional
variance budget. Its export supplies the hypotheses of the count
comparison. Thus support, the own-score and HIGH bounds, the
initializer, the resource, and the count contraction all concern one
law. An initializer or count certificate alone does not supply the
earlier hypotheses.
The additional outer-BN initializer at $n=6$ uses exactly the same
support, own-score, HIGH and predecessor premises as the ordinary
BN initializer. Its output is joined separately to the subsequent
$n=6$ resource calculation.

For $6\le n\le9$ the three bands give twelve families, each covering
the one hundred closed probability panels
$[j/100,(j+1)/100]$, $0\le j\le99$. Of these 1,200 panels,
910 are excluded by the resource inequalities, six by the separate
whole-panel sign-mass inequality of Lemma~\ref{ext:sign-mass-bound},
and 284 by complete count forests.
For $10\le n\le14$ the two bands give 1,000 panels: 798 resource
exclusions and 202 retained parent panels. The latter have 217 count
forests, because five parent panels are covered by four closed
subintervals each. These are coverage censuses, not estimates obtained
by sampling. Each retained parent rectangle is checked against the
entire corresponding count input, including its signed moment and
inward-displacement bounds. Shared endpoints belong to both adjacent
panels and cannot be discarded.

\paragraph{The complementary moment intervals.}
For $6\le n\le9$ a supported Fourier comparison covers
$[237/200,6/5]$ at radius $13/5$. Five further supported comparisons
cover the following closed intervals:
\[
\begin{array}{c|ccccc}
\rho & [6/5,7/5]&[7/5,3/2]&[3/2,8/5]&[8/5,17/10]&[17/10,7361/4000]\\
R & 3&16/5&17/5&18/5&19/5 .
\end{array}
\]
Proposition~\ref{ext:support6-14-high-table} acquires each radius by
an actual-index Gaussian calculation, using
Lemma~\ref{ext:direct-gaussian-support-criterion}.
The ensuing full Fourier calculation excludes a score above
$4097/10000<\CE$ throughout the stated band. It does not assume
the low-moment resource bound outside its domain.

For $10\le n\le14$ the supported corner is
$[23/20,71/50]$, with radius three acquired first. An unrestricted
Fourier comparison then covers $[71/50,15/8]$. Both comparisons
have target $4097/10000$.
Beyond the displayed intervals the published affine Berry--Esseen
bound already used in the entry argument gives
\[
 \frac{\sqrt n\Delta_n(P)}{\rho}
 \le\frac{3031}{10000}\left(1+\frac{323}{500\rho}\right)
 \le\frac{819}{2000}<\frac{4097}{10000}<\CE
 \quad(\rho\ge7361/4000).
\]
For the second group this tail overlaps the unrestricted Fourier
interval, since $7361/4000<15/8$. Hence every $\rho\ge1$ belongs
to a low-moment band, a complementary band, or the affine tail.
All boundaries are included. Reflection and closed-event attainment
cover the original negative-discrepancy branch as well. Every possible
first failure is therefore excluded, completing the finite induction.
\end{proof}

\paragraph{What must be checked at an interface.}
The numerical record must bind complete source and result objects, not
just a success word or a maximum value. In the original fifth- and
sixth-index chains some immutable candidate inputs contain earlier
receipt bytes. Their mathematical contents must agree with the entire
newly verified parent result, apart from explicitly identified timing
metadata. Later chains bind directly to the fresh initializer and cap
receipts. In either representation the join must preserve the actual
index, the full moment interval, both support signs, every probability
panel, and the direction of every rounded inequality. These checks are
the finite proof of applicability of the count comparison; their
meaning is independent of the operating system used to run the exact
arithmetic.
